% Vietnamese translation for OI Public Library
% Source-PDF: 其他 OTHERS/博弈论相关 - 李晓潇.pdf
\documentclass[11pt]{article}
\usepackage{oipl-vn}

\newcommand{\mex}{\operatorname{mex}}
\newcommand{\SG}{\operatorname{SG}}

\title{Một số chủ đề về lý thuyết trò chơi}
\author{Li Xiaoxiao}
\date{}

\begin{document}
\maketitle

\origtitle{Một số chủ đề về lý thuyết trò chơi}
\sourcepath{Others / Game theory related - Li Xiaoxiao}

\begin{abstract}
Tài liệu gốc là một bộ slide giới thiệu lý thuyết trò chơi trong lập trình thi
đấu. Nội dung chính gồm trò chơi tổ hợp công bằng, vị thế thắng--thua, hàm
Sprague--Grundy, tổng xor của các trò chơi con, một số ví dụ như Nim,
WordCraft và POJ 3710, cùng một phần thảo luận về cách mô hình hóa chiến lược
cho một trò chơi máy bay thời gian thực bằng rời rạc hóa và quy hoạch động.
Bản dịch chuyển slide thành ghi chú ngắn, giữ lại các định nghĩa, công thức và
ý tưởng thuật toán chính.
\end{abstract}

\tableofcontents

\section{Khởi động}

Lý thuyết trò chơi xuất hiện trong nhiều bối cảnh: kinh tế, quân sự, thi đấu,
và cả những mô hình ra quyết định trong đời sống. Trong lập trình thi đấu, ta
thường gặp các trò chơi hữu hạn giữa hai người, trong đó mọi nước đi đều biết
rõ và không có yếu tố ngẫu nhiên.

Slide gốc mở đầu bằng một ví dụ đời thường gọi là ``song đề người yêu''. Hai
người bị buộc phải chơi búa--kéo--bao để quyết định ai được sống. Họ đã hẹn sẽ
cùng ra búa, cả hai đều yêu đối phương, nhưng cả hai cũng muốn sống. Bảng điểm
minh họa trong slide là:
\[
\begin{array}{c|ccc}
 & \text{búa} & \text{kéo} & \text{bao}\\
\hline
\text{búa} &(0,0)&(1,0)&(-1,0)\\
\text{kéo}&(0,1)&(1,1)&(0,-1)\\
\text{bao}&(0,-1)&(-1,0)&(-1,-1)
\end{array}
\]
Ví dụ này không phải trọng tâm thuật toán, nhưng nhấn mạnh rằng cùng một mô
hình trò chơi có thể mô tả cả lựa chọn chiến lược lẫn bài toán thi đấu.

Một ví dụ nhỏ để khởi động là trò chơi đếm số. Hai người A và B luân phiên
nói một số nguyên dương, mỗi lần ít nhất \(1\) và nhiều nhất \(m\). Nếu tổng
các số đã nói đạt đúng \(n\) thì người vừa nói thắng. Biến thể khác quy định
người làm tổng lớn hơn hoặc bằng \(n\) sẽ thua. Các bài này cho thấy tư tưởng
cơ bản: phân loại những trạng thái mà người sắp đi có thể ép thắng và những
trạng thái mà người sắp đi không thể tránh thua.

\section{Trò chơi tổ hợp công bằng}

Một \term{trò chơi tổ hợp công bằng} có các đặc điểm sau:
\begin{itemize}
  \item hai người chơi luân phiên quyết định;
  \item từ mỗi trạng thái, tập nước đi kế tiếp là hữu hạn;
  \item các nước đi hợp lệ chỉ phụ thuộc vào trạng thái hiện tại, không phụ
        thuộc người chơi hay lịch sử ngoài trạng thái;
  \item nếu đến lượt một người mà không có nước đi hợp lệ, người đó thua;
  \item trò chơi kết thúc sau hữu hạn bước.
\end{itemize}
Điều kiện cuối cùng thường tương đương với việc đồ thị trạng thái không có
chu trình theo hướng nước đi.

\subsection{Vị thế N và P}

Ta gọi một trạng thái là \term{vị thế N} nếu người sắp đi, tức người
\emph{next}, có chiến lược thắng. Ta gọi một trạng thái là \term{vị thế P}
nếu người vừa đi trước đó, tức người \emph{previous}, có chiến lược thắng;
nói cách khác người sắp đi sẽ thua nếu cả hai chơi tối ưu.

Cách gán nhãn chuẩn là:
\begin{enumerate}
  \item mọi trạng thái kết thúc, nơi không còn nước đi, là vị thế P;
  \item nếu một trạng thái có thể đi một bước tới vị thế P, nó là vị thế N;
  \item nếu mọi nước đi từ một trạng thái đều tới vị thế N, nó là vị thế P.
\end{enumerate}
Vì trò chơi hữu hạn nên lặp lại quá trình này theo thứ tự topo sẽ gán nhãn
được mọi trạng thái.

Với trò chơi đếm số mà mỗi lần được tăng tổng thêm \(1,\ldots,m\) và ai đạt
đúng \(n\) thì thắng, các vị thế P theo khoảng cách đến \(n\) chính là các bội
của \(m+1\). Ví dụ \(m=3\), bảng đầu tiên là
\[
\begin{array}{c|ccccccccc}
n&1&2&3&4&5&6&7&8&9\\
\hline
\text{loại}&N&N&N&P&N&N&N&P&N
\end{array}
\]

\section{Nim và mô hình đồ thị}

Trong trò chơi Nim có \(n\) đống đá. Mỗi lượt, người chơi chọn một đống và
lấy đi một số đá dương tùy ý. Ai không thể đi thì thua. Với hai đống, có thể
ký hiệu trạng thái là \(f(i,j)\); với nhiều đống, điều cần phân tích là sự
khác biệt bản chất giữa các lớp trạng thái.

Một mô hình tổng quát là trò chơi trên đồ thị có hướng không chu trình. Có
một quân cờ đặt tại một đỉnh. Hai người lần lượt di chuyển quân theo một cung
có hướng; người không còn cung đi ra để di chuyển thì thua. Khi đó mỗi đỉnh là
một trạng thái của trò chơi tổ hợp công bằng.

\section{Hàm Sprague--Grundy}

Với một đồ thị có hướng không chu trình, định nghĩa hàm SG trên mỗi đỉnh
\(x\) bởi
\[
  \SG(x)=\mex\{\SG(y)\mid y \text{ là hậu duệ trực tiếp của }x\}.
\]
Ở đây \(\mex S\) là số nguyên không âm nhỏ nhất không thuộc tập \(S\).
Tương tự, với một trò chơi tổ hợp công bằng, \(\SG(x)\) của một trạng thái
\(x\) được định nghĩa bằng mex của các giá trị SG của mọi trạng thái có thể
đi tới sau một nước.

Giá trị SG biến mọi trò chơi tổ hợp công bằng hữu hạn thành một đống Nim
tương đương. Trạng thái có \(\SG(x)=0\) là vị thế P; trạng thái có
\(\SG(x)>0\) là vị thế N.

\subsection{Định lý Sprague--Grundy}

Giả sử một trò chơi tổng hợp là tổng độc lập của các trò chơi con
\[
  X=X_1+X_2+\cdots+X_n,
\]
nghĩa là mỗi lượt người chơi chọn đúng một trò chơi con và đi trong trò chơi
đó. Khi đó
\[
  \SG(X)=\SG(X_1)\oplus \SG(X_2)\oplus\cdots\oplus \SG(X_n),
\]
trong đó \(\oplus\) là xor theo bit. Vì thế hai trạng thái có cùng giá trị SG
có thể xem là tương đương về mặt trò chơi.

\section{Một số ví dụ}

\subsection{Tô màu đoạn: ZOJ 2083}

Có \(N\) đoạn thẳng với độ dài khác nhau. Hai người A và B luân phiên tô màu;
mỗi lần phải tô đúng một đoạn con dài \(2\) đơn vị và mỗi đơn vị độ dài không
được tô lại. Người không thể đi thì thua.

Khi chỉ có một đoạn, ta tính giá trị SG theo độ dài. Với nhiều đoạn, mỗi đoạn
là một trò chơi con độc lập, nên kết quả toàn cục là xor của các giá trị SG
của từng đoạn.

\subsection{RIMS}

Trên mặt phẳng có hữu hạn điểm và một số đường cong không giao nhau. Hai người
luân phiên vẽ một đường cong kín, không được cắt các đường cong đã có, và mỗi
đường cong mới phải đi qua ít nhất một điểm. Ai không thể vẽ thì thua.

Điểm đáng chú ý là kết quả chỉ phụ thuộc vào số điểm, không phụ thuộc vị trí
cụ thể. Vẽ một đường cong kín tương đương với chia tập điểm thành hai phần,
vì thế có thể tính SG theo số điểm của từng miền con.

\subsection{POJ 3710}

Bài toán cho một cây có gốc; một số nút có thể gắn với các chu trình, và các
chu trình đôi một không giao nhau. Hai người luân phiên xóa một cạnh, sau đó
bỏ toàn bộ phần không còn liên thông với gốc. Người không thể đi thì thua.

Với một dây đơn, giá trị trò chơi bằng độ dài dây:
\[
  G(X)=\text{độ dài của dây}.
\]
Với một cây, ta tách thành các cây con và lấy xor các giá trị. Với một chu
trình, xét việc xóa một cạnh trên chu trình: chu trình lẻ tương đương một dây
độ dài \(1\), còn chu trình chẵn tương đương giá trị \(0\). Do đó có thể thay
mỗi chu trình bằng một dây \(1\) hoặc \(0\), rồi tính SG của cây thu được.

\subsection{WordCraft}

Trò chơi diễn ra trên một tập xâu \(D\). Hai người luân phiên chọn một xâu
\(s\in D\), rồi xóa khỏi \(D\) mọi xâu là tiền tố của \(s\), bao gồm cả
\(s\). Người không thể đi thì thua.

Nếu dựng trực tiếp cây trò chơi, số trạng thái là quá lớn. Ngay với
\(D=\{\texttt{a},\texttt{abc},\texttt{b}\}\), ta đã phải xét nhiều tập con còn
lại sau các thao tác; với \(N\) xâu, cách nhìn thô có thể dẫn đến
\(O(2^N)\) trạng thái. Lý do là cách này chưa tận dụng từ khóa quan trọng của
bài toán: ``tiền tố''.

Cách dựng cây là mấu chốt. Mỗi xâu là một nút, và cha của nó là tiền tố dài
nhất của nó trong \(D\), khác chính nó. Khi chọn một nút, thao tác tương ứng
với xóa toàn bộ đường đi từ nút đó về gốc. Một trạng thái vì thế có thể xem
như một rừng gồm nhiều cây độc lập, và dùng định lý Sprague--Grundy để tính
kết quả.

Ví dụ trong slide dùng
\[
  D=\{\texttt{a},\texttt{aa},\texttt{ac},\texttt{acg},\texttt{acm}\}.
\]
Khi sắp các xâu theo quan hệ tiền tố, \texttt{a} là cha của \texttt{aa} và
\texttt{ac}, còn \texttt{ac} là cha của \texttt{acg} và \texttt{acm}. Một
nước đi chọn \texttt{acg} sẽ xóa đường \texttt{acg}\(\to\)\texttt{ac}\(\to\)
\texttt{a}; phần còn lại tách thành các cây con độc lập. Đây là lý do bài toán
quay về tổng xor của các trò chơi trên cây.

Thuật toán dựng cây cơ bản, với \(D[0]\) là xâu rỗng và các xâu đã sắp theo
thứ tự từ điển, có dạng:
\[
\begin{array}{l}
\textbf{FindFather}(i):\\
\quad j\gets i-1;\\
\quad \text{while }D[j]\text{ không phải tiền tố của }D[i]\text{ do}\\
\quad\quad j\gets \operatorname{father}[j];\\
\quad \operatorname{father}[i]\gets j.
\end{array}
\]
Nếu kiểm tra tiền tố trực tiếp, mỗi lần có thể mất \(O(\text{maxLen})\), và
tổng thời gian là \(O(N\text{maxLen}^2)\). Tối ưu bằng cách tính trước độ dài
tiền tố chung của \(D[i]\) và \(D[i-1]\): khi đó \(D[j]\) là tiền tố của
\(D[i]\) nếu và chỉ nếu \(|D[j]|\le \text{comLen}\). Độ phức tạp dựng cây
giảm còn \(O(N\text{maxLen})\).

Để tính SG, xử lý từ lá lên gốc. Khi xét một nút, có thể liệt kê các thao tác
bằng một lần DFS và đồng thời tính giá trị SG sau thao tác. Nếu làm thô sẽ
dẫn tới \(O(N^2)\), nhưng dùng nhận xét rằng giá trị SG của một cây con không
thể vượt quá số nút trong cây con, ta có thể giới hạn miền mex và đạt mức
hiệu quả gần \(O(N\text{maxLen})\) cho phần cốt lõi.

\section{Trường hợp trò chơi máy bay}

Phần cuối của slide chuyển sang một trò chơi thời gian thực: một bên đóng vai
\term{Boss}, bên kia đóng vai \term{Plane}. Hai vai được chơi luân phiên và
so điểm.

Ta có thể xem trò chơi thời gian thực như trò chơi theo lượt với khoảng cách
giữa hai lượt rất nhỏ. Ví dụ mỗi \(0.1\) giây là một lượt nhỏ, tối đa
\(3000\) lượt, tương đương \(5\) phút.

\subsection{Cơ chế chính}

Boss đứng yên và bắn năm loại đạn tròn từ các điểm cố định. Sức ép tăng theo
thời gian, được mô hình hóa bởi
\[
  f(t)=e^{0.00053648t},
\]
với các mốc xấp xỉ:
\[
\begin{array}{c|ccccc}
t&0&1292&2048&2585&3000\\
\hline
f(t)&1&2&3&4&5
\end{array}
\]
Mỗi lượt, số đạn mới của mỗi loại bị giới hạn bởi
\[
  \max(0.4\cdot\text{số đạn còn có thể bắn},\,6).
\]
Ví dụ nếu trên màn hình hiện có \(5\) viên loại \(1\), giới hạn hiện tại của
loại này là \(20\), thì lượt này có thể bắn tối đa
\[
  \max(0.4(20-5),6)=6
\]
viên.

Plane di chuyển thẳng đều, có vận tốc tối đa \(40\). Khi đứng yên, mỗi lượt
nó bắn một viên đạn về phía Boss. Plane nhận một điểm kỹ năng sau mỗi
\(100\) lượt. Kỹ năng tăng tốc tiêu tốn \(2\) điểm, nâng vận tốc tối đa lên
\(80\) trong \(30\) lượt; kỹ năng dọn màn hình tiêu tốn \(5\) điểm và xóa
ngay mọi viên đạn trên màn hình.

Cơ chế thắng thua của cuộc thi là hai người lần lượt đóng vai Boss và Plane,
sau đó so sánh điểm. Vì thế mục tiêu của Plane không chỉ là sống sót mà còn là
đứng yên đủ lâu để bắn và ghi điểm; mục tiêu của Boss là ép đối thủ mất điểm
hoặc bị hạ.

\subsection{Ý tưởng chiến lược}

Với Plane, chiến lược tham lam đơn giản là đứng yên để ghi điểm khi còn an
toàn; nếu nguy hiểm thì thử nhiều hướng ngẫu nhiên và chạy theo hướng thoát
được; nếu không còn đường thì dùng kỹ năng dọn màn hình hoặc chấp nhận thua.

Với Boss, hai mục tiêu tự nhiên là chặn chết Plane hoặc ép Plane phải di
chuyển liên tục để giảm điểm. Các mẫu đạn được nhắc tới gồm hình quạt dựa vào
chênh lệch vận tốc đạn, dải đạn khai thác bất đối xứng bản đồ, và các mẫu
đạn đánh lừa khiến Plane đi vào vùng dày đặc.

\subsection{Quy hoạch động cho Plane}

Một cách hệ thống hơn là rời rạc hóa mặt phẳng thành lưới \(n\times m\). Gọi
\[
  F[i][j][k]
\]
là số điểm tối đa có thể chờ thêm trong \(T\) lượt sắp tới nếu ở lượt thứ
\(k\), Plane đang tại ô \((i,j)\). Thường lấy \(T<27\) để dự báo ngắn hạn.
Giá trị cần quyết định ở trạng thái hiện tại là \(F[\text{now}_i][\text{now}_j][0]\).

Chuyển trạng thái:
\[
F[i][j][k]=
\max\begin{cases}
F[i][j][k+1]+1, & \text{nếu đứng yên tại }(i,j)\text{ an toàn},\\
F[i'][j'][k+1], & \text{nếu di chuyển tới }(i',j')\text{ hợp lệ và an toàn}.
\end{cases}
\]
Điều kiện biên là
\[
  F[i][j][T]=0.
\]

Độ phức tạp thô là
\[
  O(nmT\cdot\text{số ô trong bán kính di chuyển}\cdot\text{số đạn}).
\]
Với \(n=50\), \(m=50\), \(T=20\), số phép kiểm tra mỗi lượt có thể lên tới
hàng trăm triệu. Vì vậy cần tối ưu:
\begin{itemize}
  \item giới hạn số hướng di chuyển;
  \item dùng tìm kiếm có nhớ để loại bỏ trạng thái vô dụng;
  \item cắt tỉa bằng hàm đánh giá, cân bằng giữa tốc độ và khả năng ghi điểm;
  \item tối ưu kiểm tra va chạm với đạn, tránh kiểm tra lặp và kiểm tra không
        có ý nghĩa;
  \item đặt quy tắc hạn chế khi dùng kỹ năng.
\end{itemize}
Một vài mẹo thực tế là sau mỗi lần tránh đạn nên cố quay về vùng giữa bản đồ,
và ưu tiên ở vùng phía trên màn hình nếu điều đó giúp giảm mật độ đạn nguy
hiểm.

\section{Thông tin cuộc thi trong slide}

Các slide cuối chuyển khỏi phần thuật toán và ghi thông tin tổ chức trận chung
kết. Danh sách người tham gia được trình bày theo số thứ tự: 1 He Pufan, 2
Zhou Yichao, 3 Li Xiaoxiao, 4 Feng Qiwei, 5 Wu Lifan, 6 Fu Zhao, 7 Jiang
Xiubao, 8 Su Jiao, 9 Li Zhen, 11 Ren Yinzheng, 12 Zhou Xinyu, 13 Yan Chengji,
14 Chen Gaoyuan, 15 Lin Yuan, 16 Ren Jie. Ô số 10 trong slide nguồn để trống.
Thời gian chung kết được ghi là 7 giờ tối tại phòng báo cáo tầng một tòa nhà
chính; phần thưởng khán giả gồm Kindle Fire, máy bay điều khiển từ xa và máy
ảnh lấy liền. Slide kết thúc bằng lời cảm ơn và địa chỉ liên hệ của tác giả.

\section*{Tổng kết}

Với trò chơi tổ hợp công bằng, lối giải chuẩn là dựng đồ thị trạng thái,
phân loại vị thế N/P hoặc tính hàm Sprague--Grundy. Khi trò chơi tách thành
nhiều phần độc lập, xor các giá trị SG cho ta kết quả toàn cục. Với trò chơi
thời gian thực, cùng một tư duy mô hình hóa vẫn hữu ích: rời rạc hóa thời gian
và không gian, định nghĩa trạng thái, viết chuyển trạng thái, rồi tối ưu những
phép kiểm tra đắt nhất.

\end{document}
